% Einheit 5 von 5 – Drei Variablen und Lösungsvielfalt (Sek II) (bank/lineare-gleichungssysteme/e5.jsonl, zusammenbau v0.1)
%% TODO Zeitmarke Einheit 5: Marken-Zeile nicht lesbar
%% TODO Zweigzeile Teil 1: was hier gelernt wird (Satz in Schülersprache aus dem Einheitstitel) – Einheit 5
\einheitenkopf[e5]{Einheit 5 von 5 · Drei Variablen und Lösungsvielfalt (Sek II)}
\zweigzeile{keine Prüfungsaufgabe}
\setcounter{aufgabe}{28}

%% TODO Titel: Ich-kann-Satz fehlt in der Bank (Platzhalter: Kettenname) – Nr. 29
%% TODO Anweisung: ein Satz über der ersten Teilaufgabe fehlt in der Bank – Nr. 29
\begin{aufgabe}{Drei Variablen und Lösungsvielfalt}
\begin{teile}
\teil Wie viele Lösungen? Beim Lösen eines Systems mit drei Variablen steht am Ende $0 = 0$. Kreuze an.\\ \kreuz{genau eine}\\ \kreuz{keine}\\ \kreuz{unendlich viele}
\end{teile}
%% TODO Darstellung für schwach fehlt in der Bank (lineare-gleichungssysteme-e5-k1-s1-v1; Abschnitt „Für schwache Schüler“)
\swz[3]{Lösung nachweisen: Zeige durch Einsetzen, dass $x = 2$, $y = 1$, $z = 1$ das System I: $x + 2y - z = 3$, II: $3x + y - 2z = 5$ löst.}{}
%% TODO Darstellung für schwach fehlt in der Bank (lineare-gleichungssysteme-e5-k1-s1-v2; Abschnitt „Für schwache Schüler“)
\swz[3]{Lösung nachweisen: Zeige durch Einsetzen, dass $x = 1$, $y = 2$, $z = 3$ das System I: $x + y + z = 6$, II: $2x - y + z = 3$, III: $x + 3y - 2z = 1$ löst.}{}
%% TODO Darstellung für schwach fehlt in der Bank (lineare-gleichungssysteme-e5-k1-s1-v3; Abschnitt „Für schwache Schüler“)
\swz[3]{Lösung nachweisen: Zeige durch Einsetzen, dass $x = 1$, $y = -2$, $z = -1$ das System I: $4x - y + z = 5$, II: $x + 2z = -1$ löst.}{}
%% TODO Darstellung für schwach fehlt in der Bank (lineare-gleichungssysteme-e5-k1-s1-v4; Abschnitt „Für schwache Schüler“)
\swz[3]{Lösung nachweisen: Zeige durch Einsetzen, dass $x = 2$, $y = 2$, $z = 0$ das System I: $-x + 3y + 2z = 4$, II: $2x + y - z = 6$, III: $x - y + 3z = 0$ löst.}{}
%% TODO Darstellung für schwach fehlt in der Bank (lineare-gleichungssysteme-e5-k1-s1-v5; Abschnitt „Für schwache Schüler“)
\swz[3]{Lösung nachweisen: Zeige durch Einsetzen, dass $x = -1$, $y = 2$, $z = 2$ das System I: $x - 2y + z = -3$, II: $-x + y + 3z = 9$ löst.}{}
\swfrage{Was bleibt gleich, was ändert sich?}
%% TODO Darstellung für schwach fehlt in der Bank (lineare-gleichungssysteme-e5-k1-s2-v1; Abschnitt „Für schwache Schüler“)
\swz{$\text{Gestaffelt lösen: I: $x + y - z = 2$, II: $2y + z = 7$, III: $3z = 9$ – Lösung? Löse von unten nach oben.}$}{}
%% TODO Darstellung für schwach fehlt in der Bank (lineare-gleichungssysteme-e5-k1-s3-v1; Abschnitt „Für schwache Schüler“)
\swz{$\text{Drei Variablen: I: $2x + y = 9$, II: $y + 3z = 9$, III: $y + z = 5$ – Lösung? Subtrahiere zwei Gleichungen und setze ein.}$}{}
%% TODO Darstellung für schwach fehlt in der Bank (lineare-gleichungssysteme-e5-k1-s4-v1; Abschnitt „Für schwache Schüler“)
\swz{$\text{Wie viele Lösungen? I: $x + 2y + z = 4$, II: $x - y + 2z = 1$, III: $2x + y + 3z = 5$. Entscheide rechnerisch.}$}{}
%% TODO Darstellung für schwach fehlt in der Bank (lineare-gleichungssysteme-e5-k1-s5-v1; Abschnitt „Für schwache Schüler“)
\swz{$\text{Schar und Auswahl: I: $x - z = 1$, II: $y + z = 5$, III: $3y + 3z = 15$. Gib die Lösungsmenge mit $z = t$ an und wähle die Lösung, in der $x$, $y$ und $z$ positiv und ganzzahlig sind und $x$ am größten ist.}$}{}
%% TODO Darstellung für schwach fehlt in der Bank (lineare-gleichungssysteme-e5-k1-s6-v1; Abschnitt „Für schwache Schüler“)
\swz[3]{Erweitertes System: I: $x + y = 5$, II: $x - y = 1$ haben genau eine gemeinsame Lösung. Dazu kommt III: $2x + y = a$ mit einer reellen Zahl $a$. Wie viele Lösungen hat das System aus I, II und III, abhängig von $a$?}{}
%% TODO Darstellung für schwach fehlt in der Bank (lineare-gleichungssysteme-e5-k1-s7-v1; Abschnitt „Für schwache Schüler“)
\swz[3]{Fallunterscheidung: I: $x + y + z = 4$, II: $y - z = 2$, III: $(a^2 - 9) \cdot z = a + 3$ mit einer reellen Zahl $a$. Wie viele Lösungen, abhängig von $a$?}{}
\swz[3]{Aussage nachweisen: Eine Mischung aus drei Sorten hat die Anteile $x$, $y$ und $z$; keiner ist negativ. Das zugehörige Gleichungssystem hat die Lösungsmenge $\{(3 - 4t;\ -2 + 3t;\ t) \mid t \in \mathbb{R}\}$. Zeige: der Anteil $z$ ist mindestens doppelt so groß wie der Anteil $x$. \hfill (Abitur 2024 LK)}{}
\end{aufgabe}

%% TODO Titel: Ich-kann-Satz fehlt in der Bank (Platzhalter: Kettenname) – Nr. 30
%% TODO Anweisung: ein Satz über der ersten Teilaufgabe fehlt in der Bank – Nr. 30
\begin{aufgabe}{Drei Variablen und Lösungsvielfalt – Fehler finden, Begründen}
\swz[3]{Aus einem gestaffelten System bleibt III: $(a - 2) \cdot z = 6$. Jan rechnet \rechnung{z &= \frac{6}{a - 2}} und schreibt: „Das System hat für jedes $a$ genau eine Lösung.“ Finde den Fehler.}{}
\swz[3]{Ole löst I: $x + y + z = 2$, II: $y + 2z = 5$, III: $z = -1$: \rechnung{\text{II}: y + 2 \cdot (-1) &= 5 \\ y &= 5 - 2 = 3 \\ \text{I}: x + 3 - 1 &= 2 \\ x &= 0} Finde den Fehler.}{}
\swz[3]{Aus der Schar $x = \frac{t}{3}$, $y = 5 - t$, $z = t$ soll die Lösung mit lauter positiven ganzen Zahlen und dem größten $z$ gewählt werden. Eva nimmt $t = 4$ und erhält $x = \frac{4}{3}$, $y = 1$, $z = 4$. Finde den Fehler.}{}
\swfrage{Warum kann eine Gleichung $0 \cdot z = c$ keine oder unendlich viele Lösungen haben? Begründe.}
\swfrage{Warum liefert eine frei wählbare Variable $z = t$ unendlich viele Lösungen, und wie wählen Bedingungen wie „alle ganzzahlig und positiv“ daraus eine aus? Begründe.}
\end{aufgabe}

%% TODO Merkkasten Einheit 5: 8 Zeilen, Vorgabe höchstens fünf (3.1)
\uebersichtskasten{Drei Variablen: eine Variable aus einer Gleichung ausdrücken oder zwei Gleichungen addieren, bis ein System mit zwei Variablen übrig ist, dann wie in Einheit 2 und 3; gestaffelte Systeme (Dreiecksform) von unten nach oben durch Rückwärtseinsetzen lösen; die Lösung als Tripel (x; y; z) mit Probe in allen Gleichungen. \\ II x₂ + 2x₃ = 5, III x₂ + x₃ = 3: II − III gibt x₃ = 2, dann x₂ = 1, aus I 3x₁ − 2 = 13 folgt x₁ = 5 – Lösung (5; 1; 2). \\ Drei Fälle: führt die Rechnung auf eine wahre Aussage (0 = 0), gibt es unendlich viele Lösungen; auf eine falsche (0 = c mit c ≠ 0), keine; sonst genau eine. Zwei Gleichungen sind äquivalent, wenn eine ein Vielfaches der anderen ist – zwei nicht äquivalente Gleichungen mit zwei Variablen haben höchstens eine gemeinsame Lösung. \\ Unendlich viele Lösungen: eine Variable frei wählen (z = t), die anderen daraus ausdrücken – die Lösungsmenge ist eine Schar {(x(t); y(t); t) | t ∈ ℝ}; Zusatzbedingungen (alle negativ, ganzzahlig, größtes y) sieben aus der Schar die gesuchte Lösung. \\ I −4x + z = 4, II 2y − z = 4, III 4y − 2z = 8 (III ist das Doppelte von II): z = t, x = t/4 − 1, y = t/2 + 2; alle drei negativ und ganzzahlig heißt t < −4 und t ein Vielfaches von 4, größtes y bei t = −8: Lösung (−3; −2; −8). \\ Parameter: einen parameterabhängigen Koeffizienten nie einfach wegteilen – erst den Fall „Koeffizient null“ ansehen: (4 − a²) · z = 2 + a hat für a = 2 keine Lösung (0 = 4), für a = −2 unendlich viele (0 = 0), sonst genau eine. \\ Erweitert und verwandt: eine dritte Gleichung mit Parameter passt genau dann zur eindeutigen Lösung der ersten beiden, wenn diese sie erfüllt (t = −3 aus (1; −1) in −2x + y = t); ein Parameterwert liefert unendlich viele Lösungen, wenn die dritte Gleichung eine Kombination der beiden anderen wird (I + II gleich III für a = 5). \\ Unlösbar konstruieren: dieselbe linke Seite bis auf einen Faktor, eine andere rechte Seite – 3 · I + II* mit a = 3 und b = −3 liefert 0 = −12.}
